\abstract{
We introduce a Market Fundamentals Index (MFI) constructed from a tensor representation of firm-level accounting data and a Tucker-based smoothing step. The index summarizes time-varying dispersion in fundamentals and we test its dependence with a price-based financial chaos index (FCIX). Building on the same tensor representation, we develop a predictive modeling pipeline that imputes missing fundamentals using a mask-aware Tucker decomposition, forms quarterly lookback tensors, and forecasts next-quarter fundamentals via a low-rank CANDECOMP/PARAFAC tensor regression with firm-by-feature fixed effects and RMS normalization. The framework is evaluated against a ridge baseline using time-series cross-validation and an out-of-sample test set.
}

\keywords{multilinear tensor regression, Tucker decomposition, CANDECOMP/PARAFAC, market fundamentals, volatility index, fixed effects, time-series forecasting}

\section{Predictive Modeling for Fundamentals}
\label{Section:Predictive_Fundamentals}

We now assess whether the fundamentals tensor contains predictive structure for next-quarter accounting values. The task is to forecast a firm-by-feature target matrix using the preceding $L$ quarters of fundamentals, with $L \in \{2,4\}$ corresponding to 6-month and 1-year lookbacks.

\subsection{Tensor construction and imputation}
\label{Section:Tensor_Imputation}

We work with a quarterly panel of $N=49$ large-cap firms and $F=24$ Compustat fundamentals over 2005Q1--2024Q4. After snapping report dates to a strict quarterly grid (keeping the last report per quarter), each feature is transformed by the log-modulus map
\begin{equation}
\phi(x) = \mathrm{sgn}(x)\log(1 + |x|),
\end{equation}
which stabilizes scale while preserving sign information. The resulting raw tensor is
$\mathbfcal{X} \in \mathbb{R}^{N \times F \times T}$.
For each time index $t$, we form an input window
\begin{equation}
\mathbfcal{X}_t = \mathbfcal{X}_{:,:,t:t+L-1} \in \mathbb{R}^{N \times F \times L},
\qquad
\mathbf{Y}_t = \mathbfcal{X}_{:,:,t+L} \in \mathbb{R}^{N \times F},
\end{equation}
with observation masks $\mathbfcal{M}_t$ and $\mathbf{M}_t$ for the input and target, respectively.

Missing entries inside each window are imputed with a mask-aware Tucker decomposition,
\begin{equation}
\min_{\mathbfcal{H}_t,\mathbf{A}_t,\mathbf{B}_t,\mathbf{C}_t}
\left\| \mathbfcal{M}_t \odot \left(\mathbfcal{X}_t -
\mathbfcal{H}_t \times_1 \mathbf{A}_t \times_2 \mathbf{B}_t \times_3 \mathbf{C}_t \right) \right\|_F^2,
\end{equation}
using ranks $(40,20,L)$. We compute an observed-only RMS scale for each window,
\begin{equation}
r_t = \sqrt{\frac{1}{|\Omega_t|}\sum_{(i,j,k)\in\Omega_t} x_{ijk}^2},
\end{equation}
where $\Omega_t$ is the set of observed entries in $\mathbfcal{X}_t$. We
evaluate two input-scaling variants: an RMS-normalized covariate tensor
($\mathbfcal{X}_t / r_t$) and an unscaled covariate tensor where the imputed
values are returned to the original units.

\begin{figure}[ht]
\centering
\fbox{\rule{0pt}{2.2in}\rule{0.9\linewidth}{0pt}}
\caption[Placeholder for tensor windowing and imputation.]{Placeholder for the quarterly windowing, Tucker imputation, and RMS scaling pipeline.}
\label{Fig:Imputer_Pipeline}
\end{figure}

\subsection{Low-rank CANDECOMP/PARAFAC regression with fixed effects}
\label{Section:CP_Structured}

To control for persistent firm-feature offsets, we estimate fixed effects
\begin{equation}
\mu_{ij} = \frac{\sum_t m_{tij} y_{tij}}{\sum_t m_{tij} + \varepsilon},
\end{equation}
and model the de-meaned target with a low-rank CANDECOMP/PARAFAC tensor
regression. Let $\mathbfcal{W}$ be a low-rank coefficient tensor with CP form
\begin{equation}
\mathbfcal{W} = \sum_{r=1}^{R} \mathbf{a}_r \circ \mathbf{b}_r \circ \mathbf{c}_r \circ \mathbf{d}_r.
\end{equation}
We estimate $\mathbfcal{W}$ by regularized least squares,
\begin{equation}
\min_{\mathbfcal{W}}
\sum_t \left\| \mathbf{M}_t \odot \left(\mathbf{Y}_t - \mu - \langle \mathbfcal{X}_t, \mathbfcal{W} \rangle \right) \right\|_F^2
+ \lambda \|\mathbfcal{W}\|_F^2,
\end{equation}
where $\langle \mathbfcal{X}_t, \mathbfcal{W} \rangle$ denotes the multilinear contraction that maps the $N \times F \times L$ input to an $N \times F$ prediction. We tune the rank $R$, regularization weight $\lambda$, and optional RMS scaling via hyperparameter optimization on a rolling time-series split of the development data.

\subsection{Baseline and evaluation}
\label{Section:Baseline_Eval}

As a baseline, we fit per-feature ridge regressions on vectorized lag windows with the same firm fixed effects. Ridge penalties are selected by an inner time-series cross-validation; when a feature is too sparse, we revert to the firm-mean baseline. Performance is evaluated using a pooled, mask-aware coefficient of determination,
\begin{equation}
R^2 = 1 - \frac{\sum_{t,i,j} m_{tij} (y_{tij} - \widehat{y}_{tij})^2}
{\sum_{t,i,j} m_{tij} (y_{tij} - \bar{y})^2},
\end{equation}
with $\bar{y}$ the mean of observed entries. We also compute per-feature $R^2$ values, estimate AR(1) persistence per feature, and examine the association between persistence and the CP/PARAFAC versus ridge performance gap using Spearman correlation and paired Wilcoxon tests with Benjamini--Hochberg adjustment.

\subsection{Empirical results}
\label{Section:Predictive_Results}

Tables~\ref{Tab:CP_Ridge_Results} and~\ref{Tab:Persistence_Stats} and
Figure~\ref{Fig:Persistence_Scatter} are placeholders for the CP/PARAFAC
tensor regression versus ridge evaluation and persistence diagnostics.

\begin{table}[ht]
\centering
\caption{CP/PARAFAC tensor regression vs ridge results (placeholders).}
\label{Tab:CP_Ridge_Results}
\begin{tabular}{lcccccc}
\toprule
Input scaling & $L$ & Ridge CV $R^2$ & CP CV $R^2$ & Ridge Test $R^2$ & CP Test $R^2$ & Delta \\
\midrule
Unscaled        & 2 & -- & -- & -- & -- & -- \\
Unscaled        & 4 & -- & -- & -- & -- & -- \\
RMS-normalized  & 2 & -- & -- & -- & -- & -- \\
RMS-normalized  & 4 & -- & -- & -- & -- & -- \\
\bottomrule
\end{tabular}
\end{table}

\begin{table}[ht]
\centering
\caption{Persistence and significance diagnostics (placeholders).}
\label{Tab:Persistence_Stats}
\begin{tabular}{lcccccc}
\toprule
Input scaling & $L$ & $N_{\mathrm{feat}}$ & Mean Delta & Spearman $r$ & Spearman $p_{\mathrm{adj}}$ & Wilcoxon $p_{\mathrm{adj}}$ \\
\midrule
Unscaled        & 2 & -- & -- & -- & -- & -- \\
Unscaled        & 4 & -- & -- & -- & -- & -- \\
RMS-normalized  & 2 & -- & -- & -- & -- & -- \\
RMS-normalized  & 4 & -- & -- & -- & -- & -- \\
\bottomrule
\end{tabular}
\end{table}

\begin{figure}[ht]
\centering
\fbox{\rule{0pt}{2.2in}\rule{0.9\linewidth}{0pt}}
\caption[Placeholder for persistence vs performance scatter.]{Placeholder for persistence vs performance scatter plots (four panels: input scaling in \{unscaled, RMS-normalized\} and $L \in \{2,4\}$).}
\label{Fig:Persistence_Scatter}
\end{figure}

\section{Conclusion}

We develop a tensor-based framework for analyzing and forecasting firm fundamentals. The Market Fundamentals Index provides a fundamentals-driven view of market-wide volatility, and its statistical dependence with FCIX establishes a link between price-based and accounting-based fluctuations. We then introduce a predictive pipeline that combines mask-aware Tucker imputation with low-rank CANDECOMP/PARAFAC tensor regression and firm-by-feature fixed effects, with ridge regression as a baseline comparator. This framework enables systematic evaluation of predictability in quarterly fundamentals and can be extended to larger universes, richer feature sets, and macroeconomic covariates in future work.
